\documentclass[a4paper,11pt]{article}
\usepackage[utf8]{inputenc}
\usepackage[T1]{fontenc}
\usepackage[french]{babel}
\usepackage[left=2cm,right=2cm,top=2cm,bottom=2cm]{geometry}
\usepackage{xcolor}
\usepackage{hyperref}
\usepackage{marvosym}
\usepackage{parskip}

% --- Couleurs Valeo ---
\definecolor{primary}{RGB}{0, 120, 60}
\hypersetup{colorlinks=true, urlcolor=primary}

\begin{document}

% --- EXPÉDITEUR ---
\noindent
\begin{minipage}[t]{0.5\textwidth}
    \textbf{Loïc Aron MBASSI EWOLO}\\
    Élève-Ingénieur Bac+5 -- Double diplôme ENSEA\\[3pt]
    \Letter\ Cergy, France\\
    \Telefon\ (+33) 7 59 52 33 98\\
    \MVAt\ \href{mailto:loicaron.mbassiewolo@ensea.fr}{loicaron.mbassiewolo@ensea.fr}
\end{minipage}
\hfill
% --- DATE ---
\begin{minipage}[t]{0.4\textwidth}
    \raggedleft
    Cergy, le 11 mars 2026
\end{minipage}

\vspace{1.2cm}

% --- DESTINATAIRE ---
\hfill
\begin{minipage}[t]{0.5\textwidth}
    \raggedleft
    \textbf{Valeo -- Division Valeo Brain}\\
    Valeo Mobility Tech Center\\
    Service Recrutement\\
    Créteil, France
\end{minipage}

\vspace{1cm}

\noindent\textbf{Objet :} Candidature en alternance -- Ingénieur Développement Logiciel HMI Embarqué (dès septembre 2026)

\vspace{0.8cm}

Madame, Monsieur,

Concevoir une interface embarquée qui traduit la complexité d'un système temps réel en une expérience de conduite intuitive : c'est le défi technique qui me motive à candidater au poste d'alternant Ingénieur Développement Logiciel HMI Embarqué au sein de Valeo Brain à Créteil. En double diplôme à l'ENSEA (Systèmes Embarqués, Signal, IA) après un cursus d'ingénieur à l'École Polytechnique de Yaoundé (Mathématiques Appliquées, Génie Logiciel), je termine actuellement un stage de recherche à l'INRIA Saclay avant d'être disponible en septembre 2026 pour une alternance de 2 à 3 ans.

Mon expérience en développement embarqué et en conception d'interfaces interactives correspond directement aux missions du poste. Dans le cadre du challenge CoHoMa (robotique autonome, collaboration ENSEA/Armée française), je développe en C/C++ sur Nvidia Jetson des modules de perception temps réel intégrant SLAM, Lidar 3D et caméra de profondeur, le tout conteneurisé sous Docker pour garantir la reproductibilité. En parallèle, j'ai conçu un simulateur médical en Unity/C\# avec une interface homme-machine immersive (scénarios interactifs, retour visuel 3D temps réel), et un simulateur de flotte de véhicules en C avec visualisation OpenGL et gestion de la concurrence (mémoire partagée, signaux UNIX). Ces projets m'ont confronté aux contraintes centrales du développement HMI embarqué : latence, gestion d'événements et architecture logicielle robuste.

Au-delà de la technique, je suis familier avec les environnements agiles et internationaux. Mon stage au MILA (Institut Québécois d'IA) et ma collaboration à venir avec l'INRIA m'ont habitué à travailler en anglais au quotidien, à participer à des revues de code et à présenter des résultats techniques lors de workshops. Mon rôle de Chef de la Cellule Projets et de la Cellule IA à Polytechnique Yaoundé (plus de 50 membres) a renforcé ma capacité à coordonner des équipes pluridisciplinaires, un atout pour intégrer votre équipe en lien avec l'Égypte et l'Allemagne.

Contribuer au département Interior Experience de Valeo Brain, c'est participer à la création des interfaces de conduite de demain au sein d'un leader mondial de la mobilité intelligente. Je suis disponible dès septembre 2026 et serais ravi d'échanger lors d'un entretien sur la valeur ajoutée que je peux apporter à votre pôle Innovation.

Je vous prie d'agréer, Madame, Monsieur, l'expression de mes salutations distinguées.

\vspace{0.8cm}

\textbf{Loïc Aron MBASSI EWOLO}\\
\textit{Élève-Ingénieur ENSEA / Polytechnique Yaoundé}

\end{document}
